
\scriptsize
\begingroup
\renewcommand{\arraystretch}{1.38}
\begin{longtable}{>{\RaggedRight\arraybackslash}p{1.3cm} >{\RaggedRight\arraybackslash}p{1.5cm} >{\RaggedRight\arraybackslash}p{1.8cm} >{\RaggedRight\arraybackslash}p{1cm} >{\RaggedRight\arraybackslash}p{1.5cm} >{\RaggedRight\arraybackslash}p{2cm} >{\RaggedRight\arraybackslash}p{2cm} >{\RaggedRight\arraybackslash}p{2cm} >{\RaggedRight\arraybackslash}p{2.65cm} >{\centering\arraybackslash}p{1cm}}
\caption{Proyectos digitales, conectividad, formación e infraestructura TIC en Colombia, 2022--2026}\label{tab:proyectos-tic-colombia}\\
\toprule
\textbf{Tema} & \textbf{Proyecto} & \textbf{Entidad responsable} & \textbf{Periodo} & \textbf{Presupuesto} & \textbf{Alcance geográfico} & \textbf{Beneficiarios} & \textbf{Objetivos} & \textbf{Principales componentes} & \textbf{Fuentes} \\
\midrule
\endfirsthead
\caption[]{Proyectos digitales, conectividad, formación e infraestructura TIC en Colombia, 2022--2026 (continuación)}\\
\toprule
\textbf{Tema} & \textbf{Proyecto} & \textbf{Entidad responsable} & \textbf{Periodo} & \textbf{Presupuesto} & \textbf{Alcance geográfico} & \textbf{Beneficiarios} & \textbf{Objetivos} & \textbf{Principales componentes} & \textbf{Fuentes} \\
\midrule
\endhead
\midrule
\multicolumn{10}{r}{\emph{Continúa en la siguiente página}}\\
\endfoot
\bottomrule
\endlastfoot
Conectividad educativa & Centros Digitales & MinTIC / Fondo Único TIC & continuidad 2022–2025 & unspecified & nacional & instituciones educativas, sobre todo rurales & garantizar conectividad educativa de largo plazo & contratos 1042-2020 y 749-2022; 14.057 metas, 14.039 activas a marzo de 2025; continuidad operativa y aumento progresivo de velocidad & \parencite{mintic_informe_gestion_molina_2025} \\
Conectividad PDET & Zonas Comunitarias para la Paz & MinTIC + ART & 2022–jul. 2026 & unspecified & 161 municipios PDET & 192.000 estudiantes y docentes; impacto indirecto de 12 millones de habitantes & conectar sedes educativas y comunidades de zonas históricamente marginadas & 1.262 zonas WiFi; servicio garantizado hasta julio de 2026 & \parencite{mintic_informe_gestion_molina_2025} \\
Conectividad hogares & ConectiVIDAd para Cambiar Vidas & MinTIC / Fondo Único TIC + Internexa & 2023–2026 & COP 270.588 millones (AE2); otros componentes con montos separados & Antioquia-Urabá, Cauca, Chocó, La Guajira, Nariño, Valle del Cauca; nuevas convocatorias en Cauca, Chocó y Nariño & hogares estratos 1 y 2; comunidades organizadas; instituciones educativas públicas & ampliar Internet de banda ancha con red troncal y última milla en zonas desconectadas & estudios e ingeniería; nodos y red neutra; fibra óptica; expansión a hogares; WiFi y redes comunitarias & \parencite{mintic_informe_gestion_molina_2025} \\
Conectividad comunitaria & Juntas de Internet / Comunidades de Conectividad & MinTIC / Fondo Único TIC & 2024–2026 & COP 189.705 millones (Acuerdo 5 reportado en 2025) & 31 departamentos & hogares rurales; juntas de acción comunal, organizaciones sociales y comunidades étnicas & convertir organizaciones locales en proveedoras de su propio Internet fijo comunitario & 520 JI-CDC; meta de 78.000 hogares; infraestructura operativa y apropiación social & \parencite{mintic_informe_gestion_molina_2025} \\
Conectividad PDET / frontera & Catatumbo y Área Metropolitana de Cúcuta JI-CDC & MinTIC + convenio territorial 1211-2025 & 2025–2026 & COP 42.000 millones & 17 municipios del Catatumbo y área metropolitana de Cúcuta & 6.000 hogares; 110 instituciones educativas; 900 beneficiarios de apropiación digital & llevar conectividad satelital y comunitaria a una de las zonas más rezagadas y complejas del país & 100 comunidades de conectividad; 12 meses de servicio a hogares; 18 meses a escuelas; apropiación a juntas & \parencite{mintic_informe_gestion_molina_2025} \\
Conectividad rural comunitaria & Conectando a los no conectados & MinTIC + Unión Europea + Colnodo & 2025–2026 piloto & unspecified & Cumbal y Ricaurte en Nariño; San José del Guaviare & comunidades rurales apartadas & desplegar redes comunitarias sostenibles y apropiación digital & hasta 8 nodos por red; centro de inclusión digital; hasta 8 computadores y servidor local & \parencite{mintic_conectando_no_conectados_2025} \\
Acceso asequible & Líneas de Fomento 2.0 / 3.0 & MinTIC / Fondo Único TIC & 2024–2025 & COP 27.952 millones (2.0); 3.0 unspecified & nacional (2.0) y Ciudad Bolívar, Bogotá (3.0) & hogares vulnerables e ISP regionales con menos de 30.000 usuarios & ampliar acceso fijo residencial de banda ancha con operadores pequeños y focalización social & convocatoria PRST minoristas; interventoría; 3.0 focalizada en 4.500 hogares vulnerables de Ciudad Bolívar & \parencite{mintic_informe_gestion_molina_2025} \\
Territorial / OXI & Obras por Impuestos digitales & MinTIC + DNP + ART + contribuyentes privados & aprobados 2025 & Nariño-Cauca: COP 49.318 millones; Tame: COP 6.439 millones & Nariño, Cauca y Tame (Arauca) & 180 escuelas, 157 entidades públicas, 194 zonas WiFi, juntas de Internet; en Tame más de 1.200 beneficiarios directos & financiar infraestructuras de conectividad y apropiación en regiones apartadas vía OXI & fibra troncal, nodos de distribución y última milla, WiFi, conectividad gubernamental, talleres de apropiación & \parencite{mintic_informe_gestion_molina_2025} \\
Redes móviles & Subasta 5G y obligaciones de hacer & MinTIC + ANE + operadores (Claro, WOM, Telecall, UT Tigo-Movistar) & dic. 2023–2026 & recaudo COP 1,37 billones; obligaciones de hacer COP 389.711 millones & nacional; 241 municipios en el universo 2025 reportado para escuelas & instituciones educativas rurales, niños y niñas, usuarios de carreteras, usuarios móviles & habilitar 5G y usar obligaciones de hacer para cerrar brechas de conectividad & asignación de 320 MHz en 3.500 MHz; conexión de escuelas por 20 años; cobertura 4G en carreteras; despliegue de fibra & \parencite{mintic_5g_realidad_2023} \\
Radar & Radar 3D de Urabá y sistema asociado ADS-B & Aerocivil & 2024–2025 & unspecified & Carepa/Urabá antioqueño & navegación aérea, seguridad aérea y marítima & fortalecer vigilancia y seguridad aérea en una zona estratégica & radar 3D, torre de 30 m, cobertura primaria de 120 millas náuticas y secundaria de 250 millas; contratación 2024 previa & \parencite{aerocivil_radar_uraba_2025} \\
Infraestructura educativa digital & Computadores para Educar / Tecnologías para el Aprendizaje & Computadores para Educar & 2025 & superior a COP 170.000 millones & nacional & establecimientos educativos públicos; niños, niñas y jóvenes & ampliar dotación tecnológica e innovación educativa & convocatoria 2025; más de 14.400 equipos adquiridos para sedes; entregas directas y por sede; reutilización y robótica educativa & \parencite{computadores_para_educar_tecnologias_aprendizaje_2025} \\
IA territorial & Centros PotencIA & MinTIC + Findeter & sep. 2024–abr. 2026 & COP 234.400 millones & 60 centros en 28 departamentos & comunidades territoriales, jóvenes, emprendedores y entidades locales & construir una red nacional de inclusión digital y apropiación de tecnologías emergentes e IA & aulas de IA, laboratorios de prototipado, internet gratuito de alta velocidad, zonas colaborativas, XR, dotación y obra civil & \parencite{mintic_informe_gestion_molina_2025} \\
Formación avanzada & Talento Tech & MinTIC & 2024–jul. 2026 & COP 161.566 millones (vigencia 2025 reportada) & 9 regiones & meta mínima de 100.032 personas en 2024–2026 & reducir brecha de competencias digitales con bootcamps & programación, IA, analítica, blockchain, nube y ciberseguridad; 4 fases y 14 contratos & \parencite{mintic_informe_gestion_molina_2025} \\
Formación masiva & SENATIC & MinTIC + SENA + OIT & dic. 2023–dic. 2026 & COP 196.314 millones & nacional & meta de 303.900 personas certificadas & fortalecer competencias digitales y empleabilidad & certificación masiva; alianza con SENA y OIT; más de 144.000 personas certificadas a ago. 2025 & \parencite{mintic_informe_gestion_molina_2025} \\
Niñez y juventud & Colombia Programa & MinTIC & nov. 2023–jul. 2026 & COP 43.043 millones & nacional & 11.200 docentes y 896.000 estudiantes de colegios oficiales & llevar pensamiento computacional y programación desde preescolar a 11° & centros de referencia; micro:bit; kits Biobots y STEM; mentorías; ferias Código en Acción; red de nodos & \parencite{mintic_informe_gestion_molina_2025} \\
Formación continua & AvanzaTec & MinTIC + empresas tecnológicas vía MoU & jun. 2024–dic. 2026 & unspecified & nacional & meta de 97.000 personas certificadas & ofrecer formación corta y certificaciones con aliados tecnológicos & 45+ ofertas de formación; rutas básica-intermedia-avanzada; más de 45.000 certificaciones al corte 2025 & \parencite{mintic_informe_gestion_molina_2025} \\
IA y cuántica & ColombIA Inteligente & Minciencias & 2025–2026 & COP 24.000 millones (convocatoria 2026); 2025 unspecified en snippet & nacional, con vocación territorial & universidades, centros de I+D+i, alianzas regionales y talento de posgrado & financiar soluciones de IA y tecnologías cuánticas para retos sociales, productivos y ambientales & hasta COP 2.000 millones por proyecto; IA y cuánticas; duración máxima de 14 meses; articulación con la política nacional de IA & \parencite{minciencias_rendicion_2024_2025_pdf} \\
Seguridad digital & CiberPaz & MinTIC + Univ. de Pamplona en formaciones & 2025 & COP 8.000 millones (sensibilizaciones) + COP 6.347 millones (formaciones) & nacional & ciudadanía general; mujeres; población LGBTIQ+; ROM; indígenas; NARP; rurales; personas mayores & promover uso seguro, ético y productivo de TIC y formar en IA aplicada & talleres, festivales, hackathones, alfabetización mediática, fake news, cursos por WhatsApp y módulo virtual con enfoque en IA & \parencite{mintic_informe_gestion_molina_2025} \\
Estado digital & Servicios Ciudadanos Digitales & MinTIC + Agencia Nacional Digital & 2025 & COP 19.000 millones (2025) & nacional & ciudadanía y entidades públicas & sostener y evolucionar el ecosistema de interacción digital Estado–ciudadano & soporte, evolución, mantenimiento, operación, uso y apropiación de plataformas SCD bajo Convenio 1225-2025 & \parencite{mintic_informe_gestion_molina_2025} \\
Estado digital & Mi Colombia Digital & MinTIC & 2022–2025 & unspecified & 1.103 municipios; entidades territoriales & entidades públicas territoriales y ciudadanía & proveer sedes electrónicas y herramientas colaborativas para gobierno digital & portales institucionales; integración con trámites, SECOP, SIGEP, carpeta ciudadana y autenticación digital; 8 millones de visitas mensuales reportadas & \parencite{mintic_informe_gestion_molina_2025} \\
Data centers & Campus Bogotá de Ascenty & Ascenty & 2023–2025 & unspecified & Bogotá-Región & clientes enterprise, cloud y tecnología & ampliar infraestructura de colocation y conectividad de grado industrial & 2 data centers; 20 MW; 18.500 m²; carrier neutral; 100\% energía renovable/carbón neutral reportado en brochure & \parencite{ascenty_colombia_data_centers} \\
Data centers & BG-02 y BG-03 de Aligned/ODATA & Aligned Data Centers & 2025–2026 & unspecified & Colombia / Bogotá-Región & demanda enterprise, hyperscale y cargas de IA & expandir capacidad de data center “AI-ready” en Colombia & BG-01 existente y nuevas BG-02/BG-03 con 144 MW combinados & \parencite{aligned_colombia_data_centers} \\
\end{longtable}
\normalsize
\endgroup
